% META: {"id": "TCOMPL-ECH-CO-028", "chapitre": "TCOMPL-ECHANTILLONNAGE", "type_objet": "corrige", "exercice_ref": "TCOMPL-ECH-EX-028", "capacites_codes": ["C4"], "sources_inspiration": [], "status": "approved"}
% BEGIN-VERIFY
% from sympy import *
% x = symbols('x')
% f = x**3 - 3*x
% fp = diff(f, x)
% fpp = diff(f, x, 2)
% assert expand(fp) == 3*x**2 - 3
% assert expand(fpp) == 6*x
% assert solve(fp, x) == [-1, 1]
% assert solve(fpp, x) == [0]
% assert f.subs(x, 0) == 0
% assert f.subs(x, -1) == 2
% assert f.subs(x, 1) == -2
% END-VERIFY
\begin{corrige}{TCOMPL-ECH-EX-028}

\textbf{Question 1.} $f'(x) = 3x^2 - 3 = 3(x-1)(x+1)$ et $f''(x) = 6x$.

\commentaireMarge{$f''$ change de signe en $0$ : inflexion.}

\textbf{Question 2.} $f''(x) > 0$ sur $]0,+\infty[$ (convexe) et $f''(x) < 0$ sur $]-\infty,0[$ (concave).

Le point $x=0$ est un point d'inflexion ($f(0) = 0$).

\textbf{Question 3.} $f$ croit sur $[-1,1]$ (max local $f(1)=-2$ \ldots) — la courbe est concave puis convexe avec un point d'inflexion en $(0,0)$.

\end{corrige}
